\subsubsection{周期境界条件：ブロック循環 CCD 行列法}
\label{app:ccd_periodic}
% ═══════════════════════════════════════════════════════════════════════════════

\begin{tcolorbox}[warnbox, title={片側スキームと周期 BC の相性問題}]
§\ref{sec:ccd_bc} の片側差分境界スキームは，\textbf{非周期的な両端固定境界}（壁面・Dirichlet・Neumann）
に対して設計されている．周期境界条件（$f_0 = f_N$）に対してこのスキームをそのまま適用すると，
境界節点 $i=0,\, N$ の微分値に浮動小数点丸め起因の $\mathcal{O}(\epsilon_{\mathrm{mach}}/h)$ 程度の
誤差が残留する．

\medskip
\textbf{不安定化の機構．}
境界節点の誤差が AB2 時間積分スキームによって増幅される現象は，
Gustafsson \cite{Gustafsson1981} の境界安定性問題として知られている．
AB2 法の安定領域は純虚数軸上で $|\mathrm{Im}(\hat{\lambda})| \lesssim 0.47$ に制限されるが，
片側スキームが導入する非対称な境界行は実部正の偽固有値を生成し，AB2 の安定領域を逸脱する．
実際，表~\ref{tab:uniform_flow_comparison} に示すように，
一様流試験（$u=1$，周期 BC）でこのスキームを誤用すると
1 ステップあたり約 16 倍の誤差増幅が生じ，$t \approx 0.07$ でブローアップする．

\medskip
\textbf{根本原因：}
節点 $i=0$ の内点ステンシルが必要とする左隣り $f_{-1}$ は，
周期境界では $f_{N-1}$（折り返し）であるが，
片側スキームは代わりに $\{f_0, f_1, f_2, f_3\}$ を用いる—これは物理的に誤った構成である．
\end{tcolorbox}

\subsubsection{正しい定式化：全節点に内点スキームを適用}

周期境界条件下では，$N$ 個の独立節点（$i = 0, \ldots, N-1$；節点 $N$ は節点 0 の周期像）が
\textbf{すべて内点スキーム（式~\eqref{eq:ccd_block}）を折り返し添字で使用可能}である：
%
\begin{equation}
  \mathbf{A}_L\,\bm{v}_{(i-1)\bmod N}
  + \mathbf{B}\,\bm{v}_i
  + \mathbf{A}_R\,\bm{v}_{(i+1)\bmod N}
  = \mathbf{r}_i, \qquad i = 0, 1, \ldots, N-1
  \label{eq:ccd_periodic}
\end{equation}
%
右辺は折り返し添字で評価する：
%
\begin{equation}
  \mathbf{r}_i =
  \begin{bmatrix}
    \dfrac{a_1}{h}\bigl(f_{(i+1)\bmod N} - f_{(i-1)\bmod N}\bigr) \\[8pt]
    \dfrac{a_2}{h^2}\bigl(f_{(i-1)\bmod N} - 2f_i + f_{(i+1)\bmod N}\bigr)
  \end{bmatrix}
  \label{eq:rhs_periodic}
\end{equation}
%
求解後，周期像節点の微分値は $\bm{v}_N = \bm{v}_0$ とおく（同一物理点）．

\subsubsection{全体行列：\texorpdfstring{$2N\times 2N$}{2N x 2N} ブロック循環行列}

$N$ 式（各 $2\times 2$ ブロック）をまとめた全体行列は
\textbf{ブロック循環（block-circulant）行列}となる：
%
\begin{equation}
  \mathbf{K}_{\mathrm{circ}} =
  \begin{bmatrix}
    \mathbf{B}   & \mathbf{A}_R & \mathbf{0}   & \cdots & \mathbf{0}   & \mathbf{A}_L \\
    \mathbf{A}_L & \mathbf{B}   & \mathbf{A}_R & \cdots & \mathbf{0}   & \mathbf{0}   \\
    \mathbf{0}   & \mathbf{A}_L & \mathbf{B}   & \ddots & \vdots       & \vdots       \\
    \vdots       &              & \ddots        & \ddots & \mathbf{A}_R & \mathbf{0}   \\
    \mathbf{0}   & \cdots       & \mathbf{0}   & \mathbf{A}_L & \mathbf{B} & \mathbf{A}_R \\
    \mathbf{A}_R & \mathbf{0}   & \cdots & \mathbf{0}   & \mathbf{A}_L & \mathbf{B}
  \end{bmatrix}
  \in \mathbb{R}^{2N\times 2N}
  \label{eq:ccd_circulant}
\end{equation}
%
ここで $\mathbf{A}_L,\, \mathbf{A}_R,\, \mathbf{B}$ は式~\eqref{eq:ccd_block} の内点ブロックと完全に同一であり，
\textbf{境界専用ブロックは一切不要}となる．
ブロック三重対角行列（式~\eqref{eq:ccd_global}）との違いは隅の $\mathbf{A}_L$（右上：$i=0$ 行の折り返し）と $\mathbf{A}_R$（左下：$i=N{-}1$ 行の折り返し）の2ブロックのみである．

\begin{tcolorbox}[defbox, title={ブロック循環行列の正則性}]
$\mathbf{K}_{\mathrm{circ}}$ が正則であることを示す．

定数モード $\bm{v}_i = [c,\, 0]^\top$ を代入すると各行は
$(\mathbf{A}_L + \mathbf{B} + \mathbf{A}_R)[c,\, 0]^\top = [(1+2\alpha_1)c,\, 0]^\top$ となり，
$1+2\alpha_1 = 1 + 7/8 = 15/8 \neq 0$ だから零固有値にならない．

定数モード $\bm{v}_i = [0,\, c]^\top$ は $(1+2\beta_2) = 3/4 \neq 0$ から同様に非退化．

一般の波数 $k=1,\ldots,N-1$ の Fourier モード $\bm{v}_i = e^{2\pi\mathrm{i}ki/N}\bm{w}$ に対する固有値は，
$\hat{\mathbf{K}}_k = e^{-2\pi\mathrm{i}k/N}\mathbf{A}_L + \mathbf{B} + e^{+2\pi\mathrm{i}k/N}\mathbf{A}_R$ の固有値に対応し，
CCD スキームの修正波数解析（§\ref{sec:ccd_def}）から全波数で非零である．
ゆえに $\mathbf{K}_{\mathrm{circ}}$ は正則である．
\end{tcolorbox}

\subsubsection{数値解法：密行列 LU 分解}

$N \leq 512$ 程度では，$2N\times 2N$ 行列の密行列 LU 分解（前計算コスト $\mathcal{O}(N^3)$，
$N=32$ で $64\times64$，$N=128$ で $256\times256$）を初期化時に1回だけ実行し，
タイムステップ毎に \texttt{lu\_solve} によるバッチ後退代入（コスト $\mathcal{O}(N^2 B)$，$B$ はバッチサイズ）を適用する．

\begin{tcolorbox}[algbox, title={周期 CCD 解法の手順（実装概要）}]
\textbf{初期化（$\mathcal{O}(N^3)$，1回）：}
\begin{enumerate}
  \item $2N\times 2N$ ブロック循環行列 $\mathbf{K}_{\mathrm{circ}}$（式~\eqref{eq:ccd_circulant}）を組み立てる．
  \item $(\mathbf{L},\,\mathbf{U},\,\mathbf{P}) = \mathrm{lu}(\mathbf{K}_{\mathrm{circ}})$ を計算・保存する．
\end{enumerate}
\textbf{求解（$\mathcal{O}(N^2 B)$，タイムステップ毎）：}
\begin{enumerate}
  \item 折り返し添字で右辺 $\mathbf{R} \in \mathbb{R}^{2N\times B}$ を構築する（式~\eqref{eq:rhs_periodic}）．
  \item $\mathbf{K}_{\mathrm{circ}}\,\mathbf{V} = \mathbf{R}$ を LU 後退代入で解く．
  \item $[f'_i,\, f''_i]^\top = \mathbf{V}_{2i:2i+2}$ を取り出し，$\bm{v}_N = \bm{v}_0$ とする．
\end{enumerate}
\end{tcolorbox}

\subsubsection{検証：一様流試験}
\label{sec:verify_uniform_flow}

実装の正確性を確認するために，\textbf{一様流試験}（$u\equiv1,\; v\equiv0,\; p\equiv0$；
$32\times32$ 格子，$Re=10^6$，$We=Fr=10^{10}$，$\rho_r=\mu_r=1$，CFL$=0.3$，周期 BC）を実施した．
理論上，全ての項（対流・粘性・IPC 投影・面フラックス補正）が恒等的にゼロとなるため，
速度・圧力は全時刻で初期値のまま保持されなければならない．

\begin{table}[htbp]
\centering
\caption{一様流試験の結果比較：周期 CCD と片側スキーム誤適用}
\label{tab:uniform_flow_comparison}
\begin{tabular}{lccccc}
\toprule
スキーム & ステップ & $t$ & $\|u-1\|_\infty$ & $\|v\|_\infty$ & $\|p\|_\infty$ \\
\midrule
\multirow{3}{*}{周期 CCD（正）}
  &  1 & 0.009 & $0$              & $9\times10^{-23}$ & $0$ \\
  & 10 & 0.094 & $0$              & $9\times10^{-22}$ & $0$ \\
  & 20 & 0.188 & $0$              & $2\times10^{-21}$ & $0$ \\
\midrule
\multirow{3}{*}{片側スキーム（誤）}
  &  1 & 0.009 & $2.5\times10^{-9}$  & $2.5\times10^{-9}$  & $3.3\times10^{-8}$ \\
  &  7 & 0.066 & $2.5\times10^{-3}$  & $2.6\times10^{-3}$  & $2.8\times10^{-1}$ \\
  & 10 & 0.094 & $4.7\times10^{0}$   & $5.0\times10^{0}$   & $1.4\times10^{3}$  \\
\bottomrule
\end{tabular}
\end{table}

\noindent
周期 CCD では $\|u-1\|_\infty$ および $\|p\|_\infty$ は全ステップで浮動小数点表現の最小単位（マシンゼロ）であり，
$\|v\|_\infty$ は 20 ステップで $\mathcal{O}(10^{-21})$ と線形増加する（時間積分の丸め累積のみ）．
一方，片側スキームを誤用した場合は 1 ステップあたり約 16 倍の誤差増幅が生じ，$t\approx0.07$ でブローアップした．
